\section{Our Approach}
\label{sec:impl}

We implemented CART and C4.5 five times: \ImplSerialBase{}, \ImplSerialOpt{},
\ImplParallel{}, \ImplGpu{} and \ImplHybrid{}. All five grow exactly the same
tree. They differ in how a node's samples are kept sorted, which cuts are
scored, and where and in what order nodes are grown (Table~\ref{tab:impl}).

All implementations share the data layout and the split rules. Attribute values
are stored by attribute as 32-bit floats, and class labels are mapped to
integers. One double-precision routine computes the gain of a cut on both the
CPU and the GPU. Because the two can still differ in the last bit, gains within
$10^{-12}$ count as equal: the earlier cut wins, then the attribute with the
lower index. Every optimization below skips only cuts that provably cannot win
under this rule. Parallel builders create nodes in a nondeterministic order, so
nodes are renumbered in preorder at the end. A regression suite requires every
implementation to reproduce the same reference trees, which match scikit-learn
and Quinlan's C4.5 except where two cuts have exactly equal gain.

Pruning runs after growth, in CPU code shared by all implementations. With a
fixed $\alpha$, CART also stops early at nodes with fewer than $\alpha N$
misclassified training samples, since pruning would collapse any subtree below
them anyway. The pruned tree is the same, but those subtrees are never grown.

\begin{table}
\caption{The five implementations. \emph{Presort}: each attribute is sorted
once and kept sorted by stable partitioning. \emph{Bounds}: cuts that provably
cannot win are skipped, and every other cut is scored. \ImplHybrid{} grows
large nodes as \ImplGpu{} and small subtrees as \ImplParallel{}.}
\label{tab:impl}
\centering
\small
\setlength{\tabcolsep}{3.5pt}
\begin{tabular}{@{}llllll@{}}
\toprule
 & Sorting & Scoring & Hardware & Parallelism & Node order \\
\midrule
\ImplSerialBase{} & each node & all cuts        & 1 core      & none      & depth-first \\
\ImplSerialOpt{}  & presort   & bounds          & 1 core      & none      & depth-first \\
\ImplParallel{}   & presort   & bounds          & cores       & data, task & depth-first \\
\ImplGpu{}        & presort   & bounds          & GPU         & data      & breadth-first \\
\ImplHybrid{}     & presort   & bounds          & GPU + cores & data, task & both \\
\bottomrule
\end{tabular}
\end{table}

\subsection{\ImplSerialBase}
\label{sec:impl-simple}

\ImplSerialBase{} follows the textbook algorithm and is the baseline for the
others. At each node it sorts the node's samples by every attribute, scores
every cut, chooses the split with the shared rule, and recurses on the two
children. For a node with $n$ samples and $F$ attributes, sorting alone costs
$O(Fn\log n)$, and every comparison reads an attribute value from a random
memory location.

\subsection{\ImplSerialOpt}
\label{sec:impl-serial}

\ImplSerialOpt{} keeps the algorithm but changes how it is executed in two
ways.

\emph{Sort once.} Before growth, every attribute is presorted with a stable
radix sort. Each sorted entry holds the value together with the sample and its
class, so a sweep never has to look up a label. A node occupies the same index
range in every attribute. After a split, the winning attribute already lists
the left samples first, and every other attribute is partitioned stably to
match, which keeps both children sorted, as in SLIQ and
SPRINT~\cite{mehta1996sliq,shafer1996sprint}. An attribute that is constant on
a node is dropped for the node's entire subtree.

\emph{Skip cuts that cannot win.} The sweep over a sorted attribute skips cuts
that provably cannot beat the best cut so far. A cut between two samples of the
same class with unique values is never optimal~\cite{fayyad1992continuous}.
With two classes, the gain is convex in the class counts, so an upper bound can
rule out a whole block of consecutive cuts. Because moving one sample changes
the gain by a bounded amount, a cut far below the best also rules out the next
few cuts. Every remaining cut is scored by the exact routine. The winning sweep
also yields the class counts of both children.

Memory is allocated once: every attribute has a buffer of fixed size, and one
spare buffer lets large nodes be partitioned out of place.

\subsection{\ImplParallel}
\label{sec:impl-parallel}

\ImplParallel{} runs the \ImplSerialOpt{} builder on a pool of threads and
combines both kinds of parallelism, in the spirit of the NAP strategy of
Aldinucci et al.~\cite{aldinucci2014fastflow}. Nodes with at least 32,768
samples sweep and partition their attributes in parallel. When a node with at
least 4,096 samples is split, one child becomes a new task and the current
thread continues with the other. Data parallelism thus serves the few large
nodes near the root, and task parallelism the many small subtrees below. The
presort is parallel as well, and extra spare buffers let all threads partition
attributes at the same time.

Large nodes are memory-bound, so the builder avoids traffic between cores and
repeated reads. The thread that splits a node marks the samples going left in
its own bitmap, since one bitmap shared by all threads would put samples of
every node being split into each cache line. While a large node's attributes
are partitioned out of place, its children's sweeps read them on the way, so
each level reads a large node's data once. When a node has fewer attributes
than threads (SUSY, HIGGS, Chess), each attribute's sweep is cut into chunks
for all threads; the chunks are joined in order so that the chosen cut,
near-ties included, is exactly that of a single sweep. Threads are not pinned
to cores, and memory is not placed explicitly on NUMA nodes.

\subsection{\ImplGpu}
\label{sec:impl-gpu}

\ImplGpu{} grows the whole tree on the GPU, breadth-first: it processes all
nodes of a tree level together, so that deep levels with many small nodes still
keep the device busy. The attributes are presorted on the GPU as they are
uploaded and stored as in \ImplSerialOpt{}.

For each level, the GPU splits every node's attributes into tiles of 2,048
samples, finds the best cut of every tile, and reduces these to the best cut of
every node and attribute. Only these results go to the CPU, which applies the
shared CART or C4.5 rule and turns children that cannot be split into leaves.
The GPU then partitions the split nodes and, in the same pass, collects the
class counts for the next level. Each level therefore needs a single round
trip between CPU and GPU. As on the CPU, the sweeps skip cuts with the
boundary-point rule and, for two classes, with a bound over each whole tile,
and score every remaining cut with the exact routine. Deep levels hold mostly
tiny nodes, for which a block of 256 threads per tile would leave seven of its
eight warps idle. A level therefore numbers the tiles of nodes with at most 256
samples last, and one warp sweeps and partitions each of them; when no
attribute of a node spans more than 32 tiles, the reductions per node and
attribute also take a warp instead of a block. Data is uploaded through
page-locked staging buffers, placed on the GPU's NUMA node.

\subsection{\ImplHybrid}
\label{sec:impl-hybrid}

\ImplHybrid{} uses the same GPU builder but hands each child with fewer than
$m$ samples to the \ImplParallel{} builder, while the GPU continues with the
large nodes of the next level. A GPU level costs about the same for a thousand
mid-size nodes as for a few large ones, whereas CPU threads grow small subtrees
independently. By default, $m$ is $N/8192$ clamped to $[1024,2048]$ for CART
and $N/4096$ clamped to $[1024,4096]$ for C4.5; \ImplGpu{} is the special case
$m=1$.

Hand-offs are asynchronous. For each level, the GPU gathers the attributes that
can still split the handed-off children and copies them to host memory in the
background; once the copy is complete, CPU tasks start growing these subtrees.
The GPU does not wait for them. Nodes from both devices are renumbered in
preorder once both finish.
